\chapter*{Annotated bibliography}
\addcontentsline{toc}{chapter}{Annotated bibliography}
\markboth{Bibliography}{Bibliography}

The following classics underpin the models and methods of this monograph.

\begin{thebibliography}{99}

\bibitem{guiochon}
G.~Guiochon, A.~Felinger, D.~G.~Shirazi, A.~M.~Katti.
\textit{Fundamentals of Preparative and Nonlinear Chromatography},
2nd ed., Academic Press, 2006.
\emph{The definitive reference for the mass-balance, isotherm, and band-broadening
theory of Chapter~\ref{ch:models}.}

\bibitem{brooks}
C.~A.~Brooks, S.~M.~Cramer.
``Steric mass-action ion exchange: displacement profiles and induced salt
gradients,'' \textit{AIChE Journal} \textbf{38}(12), 1969--1978, 1992.
\emph{The original Steric-Mass-Action isotherm, Eq.~\eqref{eq:sma}.}

\bibitem{yamamoto}
S.~Yamamoto, K.~Nakanishi, R.~Matsuno.
\textit{Ion-Exchange Chromatography of Proteins}, Marcel Dekker, 1988.
\emph{Linear-gradient-elution theory and the Yamamoto method for estimating the
characteristic charge and equilibrium constant.}

\bibitem{melander}
W.~Melander, C.~Horv\'ath.
``Salt effects on hydrophobic interactions in precipitation and chromatography of
proteins,'' \textit{Archives of Biochemistry and Biophysics} \textbf{183}(1),
200--215, 1977.
\emph{The salting-out basis of the HIC law, Eq.~\eqref{eq:hic}.}

\bibitem{snyder}
L.~R.~Snyder, J.~W.~Dolan.
\textit{High-Performance Gradient Elution: The Practical Application of the
Linear-Solvent-Strength Model}, Wiley, 2007.
\emph{The linear-solvent-strength theory of reversed-phase retention,
Eq.~\eqref{eq:lss}.}

\bibitem{schmidtbox}
G.~E.~P.~Box, J.~S.~Hunter, W.~G.~Hunter.
\textit{Statistics for Experimenters: Design, Innovation, and Discovery},
2nd ed., Wiley, 2005.
\emph{Factorial designs, response surfaces, and the analysis of variance of
Section~\ref{sec:factorial}.}

\bibitem{mckay}
M.~D.~McKay, R.~J.~Beckman, W.~J.~Conover.
``A comparison of three methods for selecting values of input variables in the
analysis of output from a computer code,'' \textit{Technometrics} \textbf{21}(2),
239--245, 1979.
\emph{The introduction of Latin hypercube sampling, Section~\ref{sec:lhs}.}

\bibitem{wold}
S.~Wold, M.~Sj\"ostr\"om, L.~Eriksson.
``PLS-regression: a basic tool of chemometrics,'' \textit{Chemometrics and
Intelligent Laboratory Systems} \textbf{58}(2), 109--130, 2001.
\emph{Partial least squares, VIP scores, and cross-validated $Q^2$,
Section~\ref{sec:mv}.}

\bibitem{rasmussen}
C.~E.~Rasmussen, C.~K.~I.~Williams.
\textit{Gaussian Processes for Machine Learning}, MIT Press, 2006.
\emph{The Gaussian-process surrogate of Section~\ref{sec:bo}.}

\bibitem{jones}
D.~R.~Jones, M.~Schonlau, W.~J.~Welch.
``Efficient global optimization of expensive black-box functions,''
\textit{Journal of Global Optimization} \textbf{13}, 455--492, 1998.
\emph{Expected Improvement and the model-based optimisation loop,
Eq.~\eqref{eq:ei}.}

\bibitem{ich}
International Council for Harmonisation.
\textit{ICH Q8(R2): Pharmaceutical Development}, 2009.
\emph{The Quality-by-Design framework---design space, CPPs, CQAs, proven
acceptable ranges---that motivates the whole exercise.}

\bibitem{baranyi}
J.~Baranyi, T.~A.~Roberts.
``A dynamic approach to predicting bacterial growth in food,''
\textit{International Journal of Food Microbiology} \textbf{23}(3--4), 277--294,
1994.
\emph{The lag-phase growth model, Eqs.~\eqref{eq:baranyi-alpha}--\eqref{eq:lag}.}

\bibitem{rosso}
L.~Rosso, J.~R.~Lobry, J.~P.~Flandrois.
``An unexpected correlation between cardinal temperatures of microbial growth
highlighted by a new model,'' \textit{Journal of Theoretical Biology}
\textbf{162}(4), 447--463, 1993.
\emph{The cardinal-temperature model with inflection, Eq.~\eqref{eq:rosso}.}

\bibitem{monod}
J.~Monod.
``The growth of bacterial cultures,'' \textit{Annual Review of Microbiology}
\textbf{3}, 371--394, 1949.
\emph{The saturating substrate-limitation kinetics used in Eq.~\eqref{eq:growth}.}

\bibitem{luedeking}
R.~Luedeking, E.~L.~Piret.
``A kinetic study of the lactic acid fermentation,'' \textit{Journal of
Biochemical and Microbiological Technology and Engineering} \textbf{1}(4),
393--412, 1959.
\emph{Growth-associated plus maintenance product formation, Eq.~\eqref{eq:luedeking}.}

\bibitem{breiman}
L.~Breiman.
``Random forests,'' \textit{Machine Learning} \textbf{45}(1), 5--32, 2001.
\emph{The bagged-tree ensemble of Section~\ref{sec:ferm-analysis}.}

\bibitem{xgboost}
T.~Chen, C.~Guestrin.
``XGBoost: a scalable tree boosting system,'' \textit{Proc.\ 22nd ACM SIGKDD},
785--794, 2016.
\emph{Regularised gradient boosting, Eq.~\eqref{eq:xgb}.}

\bibitem{colbourn}
C.~J.~Colbourn.
``Combinatorial aspects of covering arrays,'' \textit{Le Matematiche}
\textbf{59}(1--2), 125--172, 2004.
\emph{Covering-array theory and bounds, Section~\ref{sec:ferm-screen}.}

\bibitem{plackett}
R.~L.~Plackett, J.~P.~Burman.
``The design of optimum multifactorial experiments,'' \textit{Biometrika}
\textbf{33}(4), 305--325, 1946.
\emph{Saturated two-level screening designs, Table~\ref{tab:ferm-screen}.}

\bibitem{cornell}
J.~A.~Cornell.
\textit{Experiments with Mixtures: Designs, Models, and the Analysis of Mixture
Data}, 3rd ed., Wiley, 2002.
\emph{Simplex designs and the Scheff\'e mixture polynomial, Eq.~\eqref{eq:scheffe}.}

\bibitem{fisher}
R.~A.~Fisher.
\textit{The Design of Experiments}, Oliver \& Boyd, 1935.
\emph{The founding text of experimental design---comparison, replication,
randomization, and the factorial principle of Section~\ref{sec:stats-doe}.}

\bibitem{mccullagh}
P.~McCullagh, J.~A.~Nelder.
\textit{Generalized Linear Models}, 2nd ed., Chapman \& Hall, 1989.
\emph{The unifying treatment of GLMs---link functions, the exponential family,
and IRLS fitting---behind Section~\ref{sec:stats-glm}.}

\bibitem{tibshirani}
R.~Tibshirani.
``Regression shrinkage and selection via the lasso,'' \textit{Journal of the
Royal Statistical Society B} \textbf{58}(1), 267--288, 1996.
\emph{The $L_1$-penalized regression of Section~\ref{sec:stats-reg} and its
sparsity / variable-selection property.}

\bibitem{esl}
T.~Hastie, R.~Tibshirani, J.~Friedman.
\textit{The Elements of Statistical Learning}, 2nd ed., Springer, 2009.
\emph{The standard reference tying together regularization, generalized linear
models, and tree ensembles as one view of the bias--variance trade-off.}

\end{thebibliography}
